\documentclass[12pt,reqno]{amsart}
\usepackage{math-review}
% Project-wide mathematical notation. No claims are certified here.

\title{Strong Compactness and Admissible Limits -- Thematic}
\author{Known-mathematics exposition benchmark}
\date{}
\begin{document}
\maketitle
\section{The question and the conditional answer}
The slogan ``estimate, then take a limit'' hides two distinct problems:
obtaining the right topology and staying in the domain of the equation.
We keep the same spaces and quantifiers throughout the comparison.

For an integer $d\geq1$, write $\rho$ for flat distance on
$\mathbb T^d=\mathbb R^d/\mathbb Z^d$ and set, for $0<\gamma<1$,
\[
 [v]_\gamma=\sup_{x\ne y}\frac{|v(x)-v(y)|}{\rho(x,y)^\gamma},\qquad
 \|u\|_{C^{2,\gamma}}
 =\max_{|\nu|\leq2}\|D^\nu u\|_\infty
   +\max_{|\nu|=2}[D^\nu u]_\gamma.
\]
Here $D^\nu$ denotes a periodic coordinate derivative of multi-index $\nu$.
Only the fixed flat structure is used below.

Matrix order is quadratic-form order, and matrix norms are Euclidean operator norms.
\begin{theorem}[Conditional sequential closedness]\label{thm:main}
Fix integers $d,r\geq1$ and $0<\beta<\alpha<1$.
On the flat torus $\mathbb T^d=\mathbb R^d/\mathbb Z^d$, let
\[
 A:C^{2,\beta}(\mathbb T^d)\longrightarrow
 C^0(\mathbb T^d;\operatorname{Sym}_r(\mathbb R))
\]
be continuous, where the target has the uniform operator norm, and set
\[
 \mathcal U=\{u\in C^{2,\beta}(\mathbb T^d):
       A(u)(x)>0\text{ for every }x\in\mathbb T^d\}.
\]
Let $\mathcal F:[0,1]\times\mathcal U\to C^0(\mathbb T^d)$ be continuous.
Suppose $t_j\in[0,1]$, $t_j\to t_\infty$, and
$u_j\in C^{2,\alpha}(\mathbb T^d)\cap\mathcal U$ satisfy, for constants
$C<\infty$ and $\delta>0$ independent of $j$,
\[
 \mathcal F(t_j,u_j)=0,\qquad
 \|u_j\|_{C^{2,\alpha}}\leq C,\qquad
 A(u_j)(x)\geq\delta I_r\quad\text{for every }j,x.
\]
Then a subsequence converges in $C^{2,\beta}$ to a function
$u_\infty\in C^{2,\alpha}\cap\mathcal U$, with
\[
 A(u_\infty)(x)\geq\delta I_r\quad\text{for every }x,
 \qquad \mathcal F(t_\infty,u_\infty)=0.
\]
The norms are the fixed periodic derivative/H\"older norms defined immediately
before this statement. This is a conditional sequential closedness assertion,
not an existence theorem or a smoothness assertion.
\end{theorem}
This theorem gives one closedness mechanism. It does not supply an openness
argument, an initial solution or the estimates needed to apply it.

\section{The topology supplied by the estimate}
For a second-order nonlinear equation, weak or uniform convergence of the
functions alone need not pass the operator. Here the required interface is explicit.
\begin{lemma}\label{lem:compactness}
On the flat torus $\mathbb T^d$ with fixed periodic derivative norms,
let $0<\beta<\alpha<1$. If a sequence of real functions satisfies
\[
 u_j\in C^{2,\alpha}(\mathbb T^d),\qquad
 \sup_j\|u_j\|_{C^{2,\alpha}}<\infty,
\]
then some subsequence converges strongly in $C^{2,\beta}(\mathbb T^d)$
to a function in $C^{2,\alpha}(\mathbb T^d)$.
The assertion does not claim convergence in the stronger $C^{2,\alpha}$ norm.
\end{lemma}
Uniform bounds and equicontinuity first give uniform limits of derivatives
through order two. Compatibility follows by integrating along coordinate segments.
For each second-derivative difference $w$, the sharper conclusion follows from
\[
 [w]_\beta\leq (2\|w\|_\infty)^{1-\beta/\alpha}
                   [w]_\alpha^{\beta/\alpha}.
\]
The first factor tends to zero and the second stays bounded. The lecture
proves this interpolation step and the extraction, rather than treating them as synonyms.

\section{Two equivalent ways to control admissibility}
A bound on the functions does not automatically bound their distance from a
matrix-cone boundary. The following two formulations do give exactly the same control.
\begin{lemma}\label{lem:margin}
Let $H_j,H\in C^0(\mathbb T^d;\operatorname{Sym}_r(\mathbb R))$,
let $\delta>0$ be fixed, and assume $H_j\to H$ uniformly in operator norm.
If $H_j(x)\geq\delta I_r$ for every $j,x$, then $H(x)\geq\delta I_r$
for every $x$. For positive definite $H_j(x)$ the same uniform lower bound
is equivalent to
\[
 \|H_j(x)^{-1}\|_{\mathrm{op}}\leq\delta^{-1}
 \quad\text{for every }j,x.
\]
\end{lemma}
For a positive definite matrix with eigenvalues $\lambda_1,\ldots,\lambda_r$,
\[
 \|H^{-1}\|_{\mathrm{op}}
 =\max_i\lambda_i^{-1}=\lambda_{\min}(H)^{-1}.
\]
Thus the lower-eigenvalue and inverse-norm formulations have identical
quantifier order: a single bound must work for every $j$ and every $x$.

They are not two different existence theories. The first is convenient for
passing quadratic-form inequalities to the limit; the second may be convenient
when an inverse matrix is already estimated by a particular equation.

\section{Synthesis: joint premises, not competing routes}
Choose the strong subsequence from Lemma~\ref{lem:compactness}. Continuity of
$A$ gives uniform matrix convergence. Testing each unit vector preserves
$A(u_\infty)\geq\delta I_r$, so the limit lies in $\mathcal U$. It follows that
\[
 \mathcal F(t_\infty,u_\infty)=\lim_j\mathcal F(t_j,u_j)=0.
\]
Compactness and cone control are both required by this proof. Calling one
``stronger'' than the other confuses different outputs: one supplies a topology,
the other supplies membership in the domain where continuity applies.

\section{What the argument does not imply}
Separate positivity is insufficient. Take $u_j=t_j=1/j$,
$A(u)=uI_r$ and $\mathcal F(t,u)=u-t$ on positive functions. The norm bound
holds, but the limit leaves the positive cone.

Nor may the convergence exponent be silently strengthened. In one periodic
coordinate, let
\[
 u_n(x)=n^{-(2+\alpha)}\cos(2\pi n x_1).
\]
Then $\|\partial_1^2u_n\|_\infty=(2\pi)^2n^{-\alpha}$ and
$[\partial_1^2u_n]_\beta=O(n^{\beta-\alpha})$. But comparing $x_1=0$
and $x_1=1/(2n)$ gives an $\alpha$-H\"older quotient
\[
 2(2\pi)^2\,2^\alpha,
\]
independent of $n$. These functions converge to zero in $C^{2,\beta}$,
not in $C^{2,\alpha}$. The theorem states the weaker topology deliberately;
it also makes no automatic claim of smoothness or convergence of the entire sequence.

\nocite{S-guide}
\bibliographystyle{plain}
\bibliography{../generated/references}
\end{document}
